%=============================================================================
% Gravity as Vacuum Electromagnetic Loading
% Revised version incorporating red-team review changes
% March 2026
%=============================================================================

\documentclass[aps,prd,twocolumn,superscriptaddress,nofootinbib]{revtex4-2}

\usepackage{amsmath,amssymb,amsfonts}
\usepackage{hyperref}
\usepackage{booktabs}
\usepackage{bm}

\newcommand{\CP}{\mathbb{C}P}
\newcommand{\Tr}{\operatorname{Tr}}
\newcommand{\Vint}{V_{\mathrm{int}}}
\newcommand{\dd}{\mathrm{d}}
\newcommand{\vE}{\mathbf{E}}
\newcommand{\vB}{\mathbf{B}}
\newcommand{\vD}{\mathbf{D}}
\newcommand{\vH}{\mathbf{H}}
\newcommand{\vJ}{\mathbf{J}}
\newcommand{\vS}{\mathbf{S}}
\newcommand{\va}{\mathbf{a}}
\newcommand{\vx}{\mathbf{x}}

\begin{document}

\title{Gravity as Vacuum Electromagnetic Loading}

\author{G.~T.~Alcock}
\affiliation{DFD Research Programme}

\date{\today}

\begin{abstract}
We develop the interpretation that gravity arises from electromagnetic
energy loading of the quantum vacuum, treated as a physical storage
medium with elastic modulus $u_0 = c^4/(8\pi G)$.  The dimensionless
scalar field $\psi$ measures the fractional loading, and the vacuum
refractive index $n = e^\psi$ follows uniquely from the multiplicative
composition law via Cauchy's exponential functional equation.  We derive
the full nonlinear DFD field equation from a stress--strain constitutive
relation of the vacuum medium, define gravitational strain and stress
in analogy with continuum mechanics, and show that the vacuum exhibits
strain stiffening at low gradients---the mechanism behind flat rotation
curves.  The vacuum stiffness $u_0 = \hbar H_0^2/(8\pi\alpha^{57}c)$
is parameter-free given the master invariant and one measured scale.
We quantify the EM--$\psi$ coupling prediction, address the vacuum
energy hierarchy within a framework consistent with the $\alpha^{57}$
suppression, and survey the 114-year intellectual lineage from Einstein's
variable-$c$ proposal through analog gravity to the present framework.
\end{abstract}

\maketitle


%=============================================================================
\section{Introduction}
\label{sec:intro}
%=============================================================================

The vacuum of quantum electrodynamics is not empty.  It has a
permittivity $\varepsilon_0$ and a permeability $\mu_0$, setting the
speed of light $c = 1/\sqrt{\varepsilon_0\mu_0}$.  In the Density
Field Dynamics (DFD) framework, these are not fixed constants of
``empty space'' but constitutive properties of a physical medium---the
quantum vacuum---whose electromagnetic response is modified by the
presence of mass-energy.

This paper develops the \emph{vacuum loading} interpretation of DFD,
in which:
\begin{enumerate}
\item Mass deposits a fractional energy loading $\psi$ in the vacuum,
modifying its refractive index to $n = e^\psi$.
\item The vacuum resists deformation with an elastic modulus
$u_0 = c^4/(8\pi G)$, the same coefficient that appears in the
Einstein field equations.
\item The nonlinear DFD field equation emerges as the mechanical
equilibrium condition for this loaded medium.
\item The vacuum stiffness is determined by the internal topology
through $u_0 = \hbar H_0^2/(8\pi\alpha^{57}c)$.
\end{enumerate}

The paper is organized as follows.  Section~\ref{sec:stiffness}
defines and derives the vacuum stiffness $u_0$.
Section~\ref{sec:constitutive} develops the constitutive relations.
Section~\ref{sec:field-eqn} derives the field equation from the
loading picture.  Section~\ref{sec:em-coupling} quantifies the
EM--$\psi$ coupling.  Section~\ref{sec:exponential} presents the
self-consistent argument for $n = e^\psi$.
Section~\ref{sec:nonlinear} treats the nonlinear regime with
stress--strain language.
Section~\ref{sec:prior-work} relates DFD to prior work.
Section~\ref{sec:discussion} discusses implications and open questions.
The dimensional consistency of all key equations is verified in
Appendix~\ref{app:dimensions}.


%=============================================================================
\section{The Vacuum Stiffness $u_0$}
\label{sec:stiffness}
%=============================================================================

\subsection{Definition}

The vacuum stiffness $u_0$ is defined as the energy density scale that
converts the dimensionless $\psi$-field gradient into physical energy
density:
\begin{equation}
\rho_\psi = u_0\,|\nabla\psi|^2, \qquad
u_0 \equiv \frac{c^4}{8\pi G}.
\label{eq:u0-def}
\end{equation}
Numerically, $u_0 \approx 4.82 \times 10^{43}$~J/m$^3$.

This is the coefficient appearing in the Einstein field equations
relating curvature to stress-energy.  In DFD, it acquires a direct
physical interpretation: $u_0$ is the \emph{elastic modulus} of the
vacuum treated as an electromagnetic storage medium.

\subsection{Expression via the master invariant}

The DFD master invariant~\cite{Alcock2025AbInitio} states:
\begin{equation}
\frac{G\hbar H_0^2}{c^5} = \alpha^{57}.
\label{eq:master-invariant}
\end{equation}
Given the master invariant and one measured scale ($H_0$), $u_0$ is
parameter-free:
\begin{equation}
u_0 = \frac{c^4}{8\pi G}
= \frac{\hbar H_0^2}{8\pi\,\alpha^{57}\,c}.
\label{eq:u0-from-invariant}
\end{equation}
This expresses $u_0$ entirely in terms of $(\hbar, c, H_0, \alpha)$
with no explicit $G$.  Since $\alpha$ is derived from the $\CP^2
\times S^3$ topology and the exponent $57 = k_{\max} - N_{\rm gen} =
60 - 3$ is topologically forced, the vacuum stiffness is ultimately a
topological quantity modulated by the single cosmological scale $H_0$.

Equivalently, in Planck units:
\begin{equation}
u_0 = \frac{\rho_P c^2}{8\pi},
\label{eq:u0-planck}
\end{equation}
where $\rho_P = c^5/(\hbar G^2)$ is the Planck density.

\subsection{The stiffness--topology chain}

The internal volume
$\Vint = \mathrm{Vol}(\CP^2)\times\mathrm{Vol}(S^3)
= 2\pi^2 \times 2\pi^2 = 4\pi^4$
enters the geometric prefactor
$K_{\rm geom} = \pi^{3/2}/24$, which determines $\alpha$ via the
spectral action on $\CP^2 \times S^3$~\cite{Alcock2025AbInitio}.  The
causal chain is:
\begin{equation}
\Vint = 4\pi^4
\;\xrightarrow{K_{\rm geom}}\; \alpha^{-1} \!=\! 137.036
\;\xrightarrow{\alpha^{57}}\; G
\;\xrightarrow{}\; u_0.
\end{equation}
The vacuum stiffness is exponentially sensitive to the internal
topology: a 1\% shift in $\alpha^{-1}$ produces a 57\% shift in $u_0$.


%=============================================================================
\section{Constitutive Relations}
\label{sec:constitutive}
%=============================================================================

In the loading picture, the vacuum is modelled as a dielectric medium
whose electromagnetic properties respond to the scalar field $\psi$.
The constitutive relations are:
\begin{equation}
\epsilon(\psi) = \epsilon_0\, n(\psi)\, e^{+\kappa\psi}, \qquad
\mu_{\rm mag}(\psi) = \mu_0\, n(\psi)\, e^{-\kappa\psi},
\label{eq:constitutive}
\end{equation}
with $n = e^\psi$.  The phase speed is preserved:
$v_{\rm ph} = 1/\sqrt{\epsilon\,\mu_{\rm mag}} = c/n$.

The parameter $\kappa$ controls the electric/magnetic split.  For
$\kappa = 0$, both sectors respond equally; for $\kappa \neq 0$, the
permittivity and permeability shift asymmetrically under loading.
The product $\epsilon\cdot\mu_{\rm mag} = \epsilon_0\mu_0\, e^{2\psi}
= n^2/c^2$ is independent of $\kappa$.

\subsection{Observables from the sector split}

The modified Coulomb field near a massive body is:
\begin{equation}
E(r) = \frac{q}{4\pi\epsilon_0\, r^2}\, e^{-(1+\kappa)\psi(r)}.
\label{eq:coulomb-modified}
\end{equation}
Near the Sun, $\psi_\odot(R_\odot) \approx 4.2\times 10^{-6}$, giving
a fractional modification $\delta E/E \approx -(1+\kappa)\times
4.2\times 10^{-6}$.  The LPI violation parameter is
$\xi(\kappa) = 1 - K_\epsilon^{(S)}\kappa + \mathcal{O}(\kappa^2)$,
where $K_\epsilon^{(S)}$ is a species-dependent sensitivity.

\subsection{The EM threshold}

The threshold for EM back-reaction on $\psi$ is determined by the
gauge--$\psi$ Lagrangian:
\begin{equation}
\eta_c = \alpha\,\sin^2\theta_W \approx \frac{\alpha}{4}
\approx 1.8\times 10^{-3},
\label{eq:eta-threshold}
\end{equation}
where $\eta \equiv U_{\rm EM}/(\rho c^2)$.  This value follows from
the electroweak mixing angle projection onto the vacuum polarization
vertex, as derived from the gauge--$\psi$ Lagrangian coupling.  We
retain the statement that $\eta_c = \alpha/4$ arises from the
gauge--$\psi$ sector without claiming an independent microscopic
derivation.


%=============================================================================
\section{The Field Equation from Vacuum Loading}
\label{sec:field-eqn}
%=============================================================================

\subsection{Linear regime}

In the Newtonian limit ($|\psi| \ll 1$, $\mu \to 1$), the vacuum
response is linear.  The ``strain'' in the medium is $|\nabla\psi|$,
and the restoring stiffness is constant: $K = c^4/(4\pi G)$.  The
resulting field equation is the Poisson equation:
\begin{equation}
\nabla^2\psi = -\frac{8\pi G}{c^2}\rho,
\end{equation}
a linear Hooke's-law response.

\subsection{Nonlinear constitutive relation}

Define the gravitational strain and stress (see
Section~\ref{sec:nonlinear} for physical interpretation):
\begin{align}
\text{Strain:} \quad & s \equiv \frac{|\nabla\psi|}{a_*},
\label{eq:strain-def} \\
\text{Stress:} \quad & \bm{\sigma} \equiv
\frac{c^4}{8\pi G}\,\mu(s)\,\nabla\psi,
\label{eq:stress-def}
\end{align}
where $a_* = 2a_0/c^2$ is the characteristic gradient scale and
$\mu(s) = s/(1+s)$ is the DFD crossover function.

The nonlinear vacuum constitutive law is:
\begin{equation}
\bm{\sigma}_{\rm grav}
= \frac{c^4}{8\pi G}\,\mu\!\left(\frac{|\nabla\psi|}{a_*}\right)
\nabla\psi.
\label{eq:constitutive-nonlinear}
\end{equation}

\subsection{Derivation of the full field equation}

The equilibrium condition
$\nabla \cdot \bm{\sigma}_{\rm grav} = -(\text{source})$
with the matter coupling
$\mathcal{L}_{\rm int} = -(c^2/2)\psi(\rho - \bar\rho)$
yields the full DFD field equation:
\begin{equation}
\nabla\cdot\!\left[\mu\!\left(\frac{|\nabla\psi|}{a_*}\right)
\nabla\psi\right]
= -\frac{8\pi G}{c^2}(\rho - \bar\rho).
\label{eq:full-field}
\end{equation}
This is the mechanical equilibrium condition for the loaded vacuum
medium: the divergence of the vacuum stress equals the applied body
force from matter.

\subsection{Energy functional}

In the Newtonian limit ($\mu \to 1$), the elastic energy density of
the loaded vacuum is:
\begin{equation}
u_\psi = \frac{1}{2}\cdot\frac{c^4}{8\pi G}\cdot|\nabla\psi|^2
= \frac{c^4}{16\pi G}|\nabla\psi|^2.
\label{eq:energy-density}
\end{equation}
This is the standard quadratic elastic energy
$u = \frac{1}{2}K s^2$, confirming $K = c^4/(8\pi G) = u_0$.


%=============================================================================
\section{Quantifying the EM--$\psi$ Coupling}
\label{sec:em-coupling}
%=============================================================================

If the vacuum is truly a storage medium, EM fields must couple to the
loading state: $\lambda \neq 1$, where $\lambda$ controls EM
back-reaction on $\psi$.  This is testable via the SQMS Q-ratio:
\begin{equation}
\frac{Q_{\rm low\text{-}f}}{Q_{\rm high\text{-}f}}\bigg|_{\rm DFD}
= 0.52 \pm 0.08,
\qquad
\frac{Q_{\rm low\text{-}f}}{Q_{\rm high\text{-}f}}\bigg|_{\rm BCS}
= 3.60.
\label{eq:Q-ratio}
\end{equation}
The sharp discrimination ($>30\sigma$ separation) arises because in
DFD with vacuum loading, the cavity's stored EM energy partially loads
the vacuum at the $\psi$-mode frequency, and the frequency dependence
of the $\psi$-channel loss inverts the Q-ratio relative to BCS theory.

Current constraints give
$|\lambda - 1| \lesssim 3 \times 10^{-5}$ (accidental bound),
with intentional search reach to $\sim 10^{-14}$.  The SQMS Q-ratio
test provides a shape measurement independent of the absolute value
of $|\lambda - 1|$.


%=============================================================================
\section{Why $n = e^\psi$: The Cauchy Functional Equation}
\label{sec:exponential}
%=============================================================================

\subsection{The composition law}

DFD postulates $n = e^\psi$ as the vacuum refractive index.  The
field equation yields $dn/d\psi = n$, whose unique solution with
$n(0) = 1$ is $n = e^\psi$.  But \emph{why} should the vacuum's
response be multiplicative rather than additive?

The answer lies in the physics of cumulative loading.  Consider two
successive loadings: first load the vacuum to state $\psi_1$,
producing index $n(\psi_1)$.  Then apply additional loading $\psi_2$
on top of the already-loaded vacuum.  The total loading is
$\psi = \psi_1 + \psi_2$ (loadings are additive in the field
variable).  The composition law
\begin{equation}
n(\psi_1 + \psi_2) = n(\psi_1) \cdot n(\psi_2)
\label{eq:composition}
\end{equation}
---sequential loading of the vacuum is multiplicative---uniquely
forces $n = e^\psi$ via Cauchy's exponential functional equation
(assuming measurability).

Specifically, the functional equation $f(x + y) = f(x)f(y)$ with
$f(0) = 1$ and $f$ continuous and monotone has the unique solution
$f(x) = e^{ax}$ for some $a \neq 0$~\cite{Kuczma2009}.  Setting
$a = 1$ by the normalization convention $\psi = \ln n$ gives
$n = e^\psi$.

The physical content is that there is no pre-existing medium with
independent atomic structure.  The vacuum \emph{is} the medium, and
$\psi$ represents the \emph{total} accumulated loading---not a
perturbation on some fixed background.  Unlike the Kerr effect in
nonlinear optics (where $n = n_0 + n_2 I$ is additive because $I$ is
a small perturbation on a pre-existing medium with baseline $n_0$),
the vacuum loading is cumulative and self-reinforcing: the medium
through which the next increment of energy propagates has already been
modified by all previous loading.

\subsection{Why the linear ansatz fails}

Suppose we tried $n = 1 + \beta\psi$ (additive, Kerr-like).  Then:
\begin{equation}
n(\psi_1 + \psi_2) = 1 + \beta(\psi_1 + \psi_2)
\neq (1 + \beta\psi_1)(1 + \beta\psi_2).
\label{eq:linear-fails}
\end{equation}
The right side expands to
$1 + \beta\psi_1 + \beta\psi_2 + \beta^2\psi_1\psi_2$, which exceeds
the left by $\beta^2\psi_1\psi_2$.  The composition law fails.

Additionally: (i) for $\psi < -1/\beta$, the refractive index becomes
negative, an unphysical regime; and (ii) the PPN parameters at second
post-Newtonian order are wrong: $n^2 = 1 + 2\beta\psi +
\beta^2\psi^2$ gives a $\psi^2$ coefficient of $\beta^2$, whereas GR
requires $2$ (twice the linear coefficient).

\subsection{Three supporting arguments}

Three independent lines of reasoning corroborate the composition law:
\begin{enumerate}
\item \textbf{Thermodynamic extensivity.}
The loading $\psi$ is an intensive thermodynamic variable.  For a
medium whose response depends only on the current state (not on the
path), the partition function factorizes over independent loadings:
$Z(\psi_1 + \psi_2) = Z(\psi_1) \cdot Z(\psi_2)$.  Since $n$ is
determined by $Z$, the multiplicative law follows.

\item \textbf{U(1) phase accumulation.}
In the microsector, $\psi$ modifies the gauge-field phase:
$\phi = \int n\,ds = \int e^\psi\,ds$.  The exponential ensures that
phase accumulated over a path is the product of contributions from
each infinitesimal segment, the unique form compatible with the group
structure of U(1) phase rotations.

\item \textbf{PPN consistency.}
With $\psi = -2\Phi/c^2$ and $n = e^\psi$, the expansion
$n^2 = e^{2\psi} = 1 + 2\psi + 2\psi^2 + \cdots$ automatically
generates the correct PPN parameters $\gamma = 1$, $\beta = 1$ with
all higher-order terms determined---no separate tuning at each PN
order.
\end{enumerate}


%=============================================================================
\section{The Nonlinear Regime: Stress--Strain Interpretation}
\label{sec:nonlinear}
%=============================================================================

The nonlinear constitutive law~\eqref{eq:constitutive-nonlinear}
admits a natural interpretation in the language of continuum mechanics.

\subsection{Definitions}

The gravitational ``strain'' is the dimensionless gradient normalized
by the saturation scale:
\begin{equation}
s = \frac{|\nabla\psi|}{a_*},
\label{eq:strain}
\end{equation}
where $a_* = 2a_0/c^2$ and $a_0 = 2\sqrt{\alpha}\,cH_0$ is the
crossover acceleration.

The gravitational ``stress'' is the vacuum's force response:
\begin{equation}
\sigma = \frac{c^4}{8\pi G}\,\mu(s)\,\nabla\psi,
\label{eq:stress}
\end{equation}
with dimensions of Pa (pressure), as verified in
Appendix~\ref{app:dimensions}.

These are physical interpretations of existing DFD quantities, not new
derivations.

\subsection{Regime structure}

The crossover function $\mu(s) = s/(1+s)$ produces two asymptotic
regimes:

\paragraph{High strain ($s \gg 1$, Newtonian regime).}
$\mu(s) \to 1$, so $\sigma \propto \nabla\psi$ (linear, Hookean).
The field equation reduces to $\nabla^2\psi = -(8\pi G/c^2)\rho$.

\paragraph{Low strain ($s \ll 1$, deep-field/MOND regime).}
$\mu(s) \approx s$, so $\sigma \propto |\nabla\psi|\nabla\psi$
(quadratic, strain-stiffening).  The field equation becomes:
\begin{equation}
\nabla\cdot\left[\frac{|\nabla\psi|}{a_*}\nabla\psi\right]
= -\frac{8\pi G}{c^2}\rho.
\end{equation}

\subsection{Strain stiffening}

The vacuum exhibits \emph{strain stiffening} at \emph{low}
gradients---the opposite of conventional nonlinear materials (which
typically soften under increasing load).  The effective stiffness per
unit strain is $\mu(s)/s = 1/(1+s)$, which diverges as $s \to 0$.

This is the physical mechanism behind flat rotation curves: at low
accelerations (low strain), the vacuum stiffens, requiring more
gradient (more acceleration) per unit source mass.  The vacuum
self-focuses the gravitational field at low accelerations, producing
flat rotation curves without dark matter.

The behavior is analogous to the Kerr self-focusing effect in
nonlinear optics ($n_2 > 0$), but inverted in the strain direction.


%=============================================================================
\section{Relation to Prior Work}
\label{sec:prior-work}
%=============================================================================

The idea that gravity might arise from or be described through the
electromagnetic properties of the vacuum has a long history.

\subsection{Historical foundations}

Einstein~\cite{Einstein1911} proposed in 1911 that the gravitational
field modifies the speed of light:
$c(\Phi) = c_0(1 + \Phi/c^2)$.  This is the earliest explicit
treatment of gravity as a variable refractive index.  DFD's
$n = e^\psi$ is the nonlinear completion of Einstein's linear formula.

Wilson~\cite{Wilson1921} proposed that mass creates a spatially
varying dielectric constant in the vacuum, a direct precursor to
DFD's constitutive-relation approach.

Gordon~\cite{Gordon1923} showed that EM waves in a moving dielectric
obey an effective ``optical metric.''  DFD inverts this: the optical
metric is fundamental, and the ``moving dielectric'' is the vacuum
$\psi$-field.

Dicke~\cite{Dicke1957} proposed that gravity could originate from
modifications to the vacuum's electromagnetic properties, suggesting
that the speed of light correlates with the gravitational potential of
the universe---the closest pre-DFD articulation of the vacuum loading
concept.

\subsection{Induced gravity and emergent approaches}

Sakharov~\cite{Sakharov1967} proposed that gravity emerges from vacuum
quantum fluctuations---``induced gravity.''  The Einstein--Hilbert
action arises as a one-loop effective action from matter fields,
providing a natural framework for DFD's claim that $G$ (and hence
$u_0$) is determined by vacuum mode structure.

Visser~\cite{Visser2002} reviewed Sakharov's programme, showing how
$G$ emerges from the spectrum of matter field fluctuations.  The
``stiffness'' in Sakharov's picture is precisely
$u_0 = c^4/(8\pi G)$.

\subsection{Polarizable vacuum and optical approaches}

de~Felice~\cite{deFelice1971} gave a systematic treatment of the
gravitational field as an optical medium with position-dependent
refractive index, deriving lensing, redshift, and Shapiro delay from
the optical analogy.  DFD elevates this from analogy to fundamental
dynamics.

Puthoff~\cite{Puthoff2002} developed a ``polarizable vacuum'' model
in which mass modifies the vacuum's permittivity and permeability via
a scalar function, reproducing linearized GR\@.  This is formally
similar to DFD's constitutive relations, though DFD provides a
dynamical action principle and microsector derivation.

Perlick~\cite{Perlick2000} established the rigorous mathematical
foundation for ray optics in curved spacetime via Fermat's principle.

\subsection{Analog gravity and vacuum energy}

Volovik~\cite{Volovik2003} demonstrated that superfluid $^3$He-A
provides an exact analog of gravity, electromagnetism, and the
Standard Model as collective modes of the quantum vacuum.  Crucially,
Volovik showed that the vacuum energy computed from mode sums is
cancelled by the equilibrium condition---directly supporting DFD's
distinction between $u_0$ (elastic modulus) and $\rho_\Lambda$
(residual strain), and anticipating the $\alpha^{57}$ suppression
mechanism.

Barcel\'{o}, Liberati, and Visser~\cite{Barcelo2011} provided a
comprehensive review of analog gravity models.  Electromagnetic
analogs (dielectric media with gradient refractive index) are
discussed as gravitational simulators.  DFD makes the inverse claim:
the gravitational $\psi$-field \emph{is} the physical refractive
medium, not merely an analogy.

\subsection{Spectral geometry}

Connes and Chamseddine~\cite{Connes1997,Chamseddine2008} developed the
spectral action framework from which DFD's $\alpha$
derivation~\cite{Alcock2025AbInitio} descends.  The spectral action on
a product geometry $M^4 \times F$ simultaneously produces the
Einstein--Hilbert gravitational action and the Standard Model gauge
kinetic terms.  DFD identifies $F$ with $\CP^2 \times S^3$ and uses
specific heat-kernel coefficients to derive $\alpha^{-1} = 137.036$.

Recent work by Garc\'{i}a-Compean et al.~\cite{Garcia2024} and
De~Rham et al.~\cite{DeRham2025} extends the optical-gravity analogy
into the metamaterial regime.


%=============================================================================
\section{Discussion}
\label{sec:discussion}
%=============================================================================

\subsection{The vacuum energy hierarchy}

The vacuum loading picture distinguishes three energy density scales:
\begin{center}
\begin{tabular}{lcc}
\toprule
\textbf{Scale} & \textbf{Value (J/m$^3$)} & \textbf{Role} \\
\midrule
QFT cutoff $\rho_P c^2$ & $\sim 10^{113}$ & Naive mode-sum \\
Stiffness $u_0$ & $\sim 10^{43}$ & Elastic modulus \\
Dark energy $\rho_\Lambda c^2$ & $\sim 10^{-9}$ & Residual strain \\
\bottomrule
\end{tabular}
\end{center}

The critical distinction: $u_0$ is the elastic modulus (resistance to
deformation), not the equilibrium energy density.  The observed
$\rho_\Lambda \sim 10^{-9}$~J/m$^3$ is the residual loading, a
``strain'' of order $H_0^2/c^2 \sim 10^{-52}$~m$^{-2}$.  The vacuum
loading picture provides a natural framework for understanding the
vacuum energy hierarchy, consistent with the $\alpha^{57}$ suppression:
the physical microsector has only $k_{\max} = 60$ modes, and the
residual vacuum energy is suppressed by $\alpha^{57} \sim 10^{-122}$
relative to the Planck scale, rather than being set by a continuum
mode sum up to $M_P$.

This parallels Volovik's argument~\cite{Volovik2003} that in
superfluid $^3$He, the vacuum energy from mode sums vastly
overestimates the gravitational effect, because most of the vacuum
energy is cancelled by the equilibrium condition.

\subsection{The vacuum stiffness and the master invariant}

Given the master invariant $G\cdot\hbar\cdot H_0^2/c^5 =
\alpha^{57}$~\cite{Alcock2025AbInitio} and one measured scale, $u_0$
is parameter-free:
$u_0 = \hbar H_0^2/(8\pi\alpha^{57}c)$.  Whether this expression
admits an independent derivation from the microsector topology remains
open.

\subsection{The chain of reasoning}

The vacuum loading interpretation can be stated as a single chain:
\begin{enumerate}
\item The vacuum is a storage medium with $\varepsilon_0$, $\mu_0$.
\item Mass loads the medium: $\psi = \ln n$.
\item The composition law forces $n = e^\psi$ (Cauchy).
\item The stiffness is topological:
$u_0 = \hbar H_0^2/(8\pi\alpha^{57}c)$.
\item The cosmological constant is residual strain.
\item QFT overestimates by overcounting modes.
\item EM--$\psi$ coupling is predicted and testable.
\end{enumerate}

\subsection{Experimental programme}

The key predictions are:
\begin{itemize}
\item SQMS Q-ratio: $0.52 \pm 0.08$ (DFD) vs.\ $3.60$ (BCS).
\item LPI slope $\xi(\kappa)$: species-dependent if $\kappa \neq 0$.
\item TE/TM cavity asymmetry: $\sim 2\kappa\psi$ frequency splitting.
\item Polarization rotation through $\nabla\psi$ gradients.
\end{itemize}

\subsection{Open questions}

\begin{enumerate}
\item The exact value of $\lambda - 1$ requires computing the
EM--$\psi$ mode overlap on $\CP^2 \times S^3$.
\item The sector-split parameter $\kappa$ may be computable from the
constitutive relations on the internal space.
\item Time dependence of $u_0$: if $G(t)$ varies with cosmic epoch,
implications for nucleosynthesis need quantification.
\item Whether the vacuum loading picture and Sakharov's programme can
be unified through a single spectral action calculation.
\end{enumerate}


%=============================================================================
\section{Conclusions}
\label{sec:conclusions}
%=============================================================================

We have developed the vacuum loading interpretation of DFD, in which
the gravitational $\psi$-field arises from fractional electromagnetic
energy loading of the quantum vacuum.  The exponential refractive law
$n = e^\psi$ follows uniquely from the multiplicative composition law
via Cauchy's functional equation, supported by three independent
arguments (thermodynamic extensivity, U(1) phase, PPN consistency).
The vacuum stiffness $u_0 = \hbar H_0^2/(8\pi\alpha^{57}c)$ is
determined by topology plus one scale measurement.  The nonlinear
regime admits a continuum-mechanics interpretation in which the vacuum
exhibits strain stiffening at low gradients---the mechanism behind
flat rotation curves.  The framework provides a natural understanding
of the vacuum energy hierarchy consistent with the $\alpha^{57}$
suppression, and makes falsifiable predictions for the SQMS Q-ratio
test and cavity--atom comparisons.


%=============================================================================
\begin{acknowledgments}
DFD Research Programme.
\end{acknowledgments}


%=============================================================================
\appendix
\section{Dimensional Consistency of Key Equations}
\label{app:dimensions}
%=============================================================================

We verify the SI dimensional consistency of every key equation.

\begin{center}
\begin{tabular}{lp{4.2cm}c}
\toprule
\textbf{Equation} & \textbf{Dimensions} & \textbf{Status} \\
\midrule
$u_0 = c^4/(8\pi G)$ &
$\mathrm{m^4 s^{-4}}/\mathrm{(m^3 kg^{-1} s^{-2})}$
$= \mathrm{kg\,m^{-1}\,s^{-2}}$
$= \mathrm{J/m^3}$ & \checkmark \\
\addlinespace
Strain $s = |\nabla\psi|/a_*$ &
$\mathrm{m^{-1}/m^{-1}}$ = dimensionless & \checkmark \\
\addlinespace
Stress $\sigma = \frac{c^4}{8\pi G}\mu\nabla\psi$ &
$\mathrm{J/m^3}\cdot\mathrm{m^{-1}}$
$= \mathrm{N/m^2}$
$= \mathrm{Pa}$ & \checkmark \\
\addlinespace
Field eq.\ LHS:
$\nabla\!\cdot\!(\mu\nabla\psi)$ &
$\mathrm{m^{-1}\cdot m^{-1}} = \mathrm{m^{-2}}$ & \checkmark \\
\addlinespace
Field eq.\ RHS:
$8\pi G\rho/c^2$ &
$\mathrm{m^3 kg^{-1} s^{-2}\cdot kg\,m^{-3}}$
$/\mathrm{m^2 s^{-2}}$
$= \mathrm{m^{-2}}$ & \checkmark \\
\addlinespace
Energy density $u_\psi$ &
$\mathrm{J/m^3}\cdot 1$
$= \mathrm{J/m^3}$ & \checkmark \\
\addlinespace
Master invariant
$G\hbar H_0^2/c^5$ &
$\frac{\mathrm{m^3 kg^{-1} s^{-2}\cdot
kg\,m^2\,s^{-1}\cdot s^{-2}}}
{\mathrm{m^5 s^{-5}}}$
$= 1$ & \checkmark \\
\addlinespace
$a_0 = 2\sqrt{\alpha}\,cH_0$ &
$\mathrm{m\,s^{-1}\cdot s^{-1}} = \mathrm{m\,s^{-2}}$ & \checkmark \\
\addlinespace
$a_* = 2a_0/c^2$ &
$\mathrm{m\,s^{-2}/m^2\,s^{-2}} = \mathrm{m^{-1}}$ & \checkmark \\
\addlinespace
$\eta_c = \alpha\sin^2\theta_W$ &
dimensionless & \checkmark \\
\addlinespace
Coulomb modified &
$\mathrm{V/m}$ (exponential is dimensionless) & \checkmark \\
\addlinespace
LPI parameter $\xi$ &
dimensionless & \checkmark \\
\addlinespace
$k_a = 3/(8\alpha)$ &
dimensionless & \checkmark \\
\addlinespace
$k_a a_0^2 = \tfrac{3}{2}(cH_0)^2$ &
$\mathrm{m^2 s^{-4}}$ both sides & \checkmark \\
\bottomrule
\end{tabular}
\end{center}

All equations are dimensionally consistent in SI units.


%=============================================================================
\begin{thebibliography}{99}

\bibitem{Einstein1911}
A.~Einstein,
``\"{U}ber den Einfluss der Schwerkraft auf die Ausbreitung des Lichtes,''
Ann.\ Phys.\ (Leipzig) \textbf{35}, 898 (1911).

\bibitem{Wilson1921}
H.~A.~Wilson,
``An electromagnetic theory of gravitation,''
Phys.\ Rev.\ \textbf{17}, 54 (1921).

\bibitem{Gordon1923}
W.~Gordon,
``Zur Lichtfortpflanzung nach der Relativit\"{a}tstheorie,''
Ann.\ Phys.\ (Leipzig) \textbf{72}, 421 (1923).

\bibitem{Dicke1957}
R.~H.~Dicke,
``Gravitation without a principle of equivalence,''
Rev.\ Mod.\ Phys.\ \textbf{29}, 363 (1957).

\bibitem{Sakharov1967}
A.~D.~Sakharov,
``Vacuum quantum fluctuations in curved space and the theory of
gravitation,''
Dokl.\ Akad.\ Nauk SSSR \textbf{177}, 70 (1967)
[Sov.\ Phys.\ Dokl.\ \textbf{12}, 1040 (1968)].

\bibitem{deFelice1971}
F.~de~Felice,
``On the gravitational field acting as an optical medium,''
Gen.\ Relativ.\ Gravit.\ \textbf{2}, 347 (1971).

\bibitem{Connes1997}
A.~H.~Chamseddine and A.~Connes,
``The spectral action principle,''
Commun.\ Math.\ Phys.\ \textbf{186}, 731 (1997).

\bibitem{Perlick2000}
V.~Perlick,
\emph{Ray Optics, Fermat's Principle, and Applications to General
Relativity},
Lecture Notes in Physics Vol.~61 (Springer, Berlin, 2000).

\bibitem{Puthoff2002}
H.~E.~Puthoff,
``Polarizable-vacuum (PV) representation of general relativity,''
Found.\ Phys.\ \textbf{32}, 927 (2002);
arXiv:gr-qc/9909037.

\bibitem{Visser2002}
M.~Visser,
``Sakharov's induced gravity: A modern perspective,''
Mod.\ Phys.\ Lett.\ A \textbf{17}, 977 (2002);
arXiv:gr-qc/0204062.

\bibitem{Volovik2003}
G.~E.~Volovik,
\emph{The Universe in a Helium Droplet}
(Oxford University Press, Oxford, 2003).

\bibitem{Chamseddine2008}
A.~H.~Chamseddine and A.~Connes,
``Why the Standard Model,''
J.\ Geom.\ Phys.\ \textbf{58}, 38 (2008).

\bibitem{Barcelo2011}
C.~Barcel\'{o}, S.~Liberati, and M.~Visser,
``Analogue Gravity,''
Living Rev.\ Relativ.\ \textbf{14}, 3 (2011).

\bibitem{Garcia2024}
H.~Garc\'{i}a-Compean et al.,
``Optical refractive index and spacetime geometry,''
arXiv:2408.10723 (2024).

\bibitem{DeRham2025}
C.~De~Rham et al.,
``Gravitational metamaterials from optical properties of spacetime media,''
arXiv:2504.09987 (2025).

\bibitem{Alcock2025AbInitio}
G.~T.~Alcock,
``Ab initio derivation of the fine-structure constant from density
field dynamics,''
DFD Preprint (2025).

\bibitem{Kuczma2009}
M.~Kuczma,
\emph{An Introduction to the Theory of Functional Equations and
Inequalities}, 2nd ed.\ (Birkh\"{a}user, Basel, 2009).

\end{thebibliography}

\end{document}
